\section{Open problems}\label{sec:open}

\Cref{thm:T} settles depth reduction for every bounded number of walks and every bounded number of stacks. The questions it
leaves are quantitative, algorithmic, or about other layouts; the first item is (T) itself.

\begin{enumerate}[label=\arabic*.,leftmargin=2em]
\item \emph{Decide (T).} Construct a machine that transposes $d\times d$ matrices with $b(d)$-bit entries in time $O(L)$, or prove
that none exists. The first paper's frozen checker accepts either answer~\cite{FirstPaper2026}. By \cref{sec:consequences},
depth reduction is no longer an obstacle on the way to a lower bound for oblivious machines, at the level of schedules or of
codes (\cref{cor:HV}): what remains is the coding side.
To exclude the oblivious machines with $k$ work tapes, value width at most~$\beta$ and running time $C(L+1)$, it suffices to
prove $\Crux(\mathcal S_{k,C,\beta,D},\rho)$ for a single $\eps\in(0,1]$, at the depth~$D$ of \cref{cor:oblivious} and some rate
$\rho>\eps$ (\cref{rem:rate}), and to exclude all oblivious linear-time transposers, to do so for all $k$, $\beta$ and~$C$. The
native half bound would give the crux at rate~$\frac12$ at every depth (\cref{rem:native}); for machines with a linear data path
the linear crux suffices, which the eighth paper states as a question of relative rigidity~\cite[Problem~10.6]{EighthPaper2026}.
Conversely, an oblivious linear-time transposer must carry, for every~$\eps$, codes of rate at least $1-\eps$ on its own cone
relations at that depth (\cref{sec:cons-construction}). Adaptive machines are reached only through single traces
(\cref{rem:native}).

\item \emph{The order of the depth.} With the parameters of \cref{def:params}, the depth $D(k,\eps)=K^P-1$ satisfies
$\log\log D(k,\eps)=O(k^2\log(k/\eps))$ for $\eps\le\frac12$: the number of digits is $P=(O(k^2/\eps))^{4k^2+1}$, since
$Q=O(k^2/\eps)$ and $j=O(Q^{4k^2}\log(k/\eps))$, and the radix satisfies $\log K=O(k^2/\eps+k^2\log(k/\eps))$.
\begin{problem}[The true depth]\label{prob:depth}
What is the smallest $D(k,\eps)$ for which $\TCP(k)$ holds at~$\eps$? Is it bounded by $2^{\poly(k/\eps)}$, or even by
$\poly(k/\eps)$?
\end{problem}
The doubly exponential bound comes from Lemma~D (\cref{lem:D}), whose tail decays at the rate $1-Q^{-n^2}$ per block of $n$
steps: a block stops the front process only when all of its at most $n^2$ fresh bits are broken. A rate of $1-Q^{-O(n)}$ per
block would lower the exponent $4k^2$ in~$P$ to $O(k)$, and so give $\log\log D(k,\eps)=O(k\log(k/\eps))$; a rate bounded away
from~$1$ would give $P=O(k^2\log(k/\eps))$ and, since $\log K=O(k^2/\eps+\log P)$, a depth $2^{\poly(k/\eps)}$. Heuristically, a
better rate is plausible, since a front process advances only while some head has an unbroken candidate, and each head offers
at most one candidate per value.

\item \emph{Other layouts.} \Cref{thm:T} covers DAGs with a spine whose arcs form a bounded number of stacks. Bounded queue
number is the natural next class: expanders with bounded queue number exist~\cite{DSW2016}, but expansion is not
depth-robustness in the sense of TCP, since the spine is free and only arcs are counted.
\begin{problem}[Queues]\label{prob:queues}
Is every family of DAGs with a spine whose arcs form a bounded number of queues depth-reducible?
\end{problem}

\item \emph{Explicit cuts.} The proof is an averaging argument over the random choices $(r,\omega,\xi)$ (\cref{prop:avg}), and
the label of one choice is easy to evaluate. Finding a good choice deterministically, and in near-linear time, is open; the
number of choices is $h_0K^PQ^{\iota}$, where $\iota$, the number of pairs of a head and one of its frames, grows with~$T$.
\end{enumerate}
